\section{求解器实现与完整求解过程}\label{sec:supp-solution-zh}

本节先区分三种角色，再把“DAG 如何变成可验证工件”逐步展开。三种角色是：
默认使用的可扩展组合、只用于机制研究的离线代价修复，以及固定拓扑序的
CP-SAT oracle。后两者都不是生产 \texttt{solve} 的隐藏阶段。

\subsection{Canonical 可扩展求解器}

生产实现为 P2/P3 生成四个结构拓扑序：\texttt{unlock-frontier}、
\texttt{capfit-id}、P1 memory order 和 \texttt{id-raw}。只有第一个是 v2
新增排序；其余三个提供历史与结构多样性，不能重新命名为代价感知排序。
四个排序均不读取 backed/generated 标签。类别信息只进入两类 victim 评分
以及完整候选的真实流量计算与最终选择。这里的“纯结构”限定于新增的
\texttt{unlock-frontier}：portfolio 保留的旧 P1 fallback 对特定输入形态
含 \texttt{D2S}、\texttt{COPY-IN} 和 \texttt{MOVE} 相关 tie-break。因此
本文不把整个 canonical solver 称为 motif-agnostic；能成立的窄结论只是
frontier 不读取实例名、操作模式或缓冲区类别。

\texttt{unlock-frontier} 先调度就绪 FREE，再调度最小编号的就绪操作；仅
剩 ALLOC 时，寻找编号最小且“全部尚未完成前驱都已经就绪”的消费者，并
完成其前驱组。这样可避免连续的一输入搬运长期饿死一个多输入操作。

\begin{algorithm}[H]
\caption{依赖前沿拓扑序}\label{alg:supp-frontier-zh}
\begin{algorithmic}[1]
\Require DAG 实例 $I$
\State 将零入度节点划分为 $Q_f,Q_o,Q_a$；$S\gets[\ ]$
\While{$Q_f\cup Q_o\cup Q_a\ne\emptyset$}
  \If{$Q_f\ne\emptyset$} \State $v\gets Q_f$ 中最大编号节点
  \ElsIf{$Q_o\ne\emptyset$} \State $v\gets Q_o$ 中最小编号节点
  \Else
    \State $u\gets$ 尚未完成前驱全在 $Q_a$ 中的最小编号消费者
    \State $v\gets u$ 的最小编号就绪前驱
    \If{不存在这样的 $u$} \State $v\gets Q_a$ 中最小编号节点 \EndIf
  \EndIf
  \State 将 $v$ 追加至 $S$；更新入度及 $Q_f,Q_o,Q_a$
\EndWhile
\State \Return $S$
\end{algorithmic}
\end{algorithm}

每个拓扑序分别运行真正的 best-fit 地址放置：从所有足够大的空闲区间中
选择长度最小者，长度相同时选择地址最小者；释放后合并相邻空闲区间。若
没有连续空洞，放置器从当前不被下一操作 pinned 的缓冲区中选择 victim，
插入 spill-out，释放其整段地址，直至请求可放置。需要再次访问或 FREE 时
插入 spill-in，并允许换入到新地址。

生产组合比较 \texttt{dist-size-cost} 与 \texttt{share-adaptive-25} 两种
victim 策略。前者联合下次使用距离、大小和静态流量代价；后者在有备份
字节稀少或纯生成缓冲区尺寸高度异质时切换优先级。P2 枚举提前重载窗口
$H\in\{0,5,40\}$，P3 枚举 $H\in\{0,5,40,80,120\}$。提前重载只有在
已有足够连续空洞时才发生，它本身不驱逐其他驻留缓冲区；但它会改变之后的
驻留状态，因而仍须把每个窗口作为完整候选重新评价。

\begin{algorithm}[H]
\caption{P2/P3 有界工件组合}\label{alg:supp-portfolio-zh}
\begin{algorithmic}[1]
\Require 实例 $I$，问题编号 $q\in\{2,3\}$
\State $\mathcal O\gets\{$frontier, capfit-id, P1, id-raw$\}$
\State $\mathcal P\gets\{$dist-size-cost, share-adaptive-25$\}$
\State $\mathcal H\gets\{0,5,40\}$；若 $q=3$，加入 $80,120$
\State $A^\star\gets\varnothing$
\For{$(O,P,H)\in\mathcal O\times\mathcal P\times\mathcal H$}
  \State 生成拓扑序 $\pi_O$
  \State 用 best-fit、策略 $P$ 和窗口 $H$ 生成地址与溢出工件 $A$
  \If{放置失败} \State 跳过该不可行组合成员 \EndIf
  \State 计算 $E,n_{\mathrm{spill}},T$（候选循环内部不调用 canonical evaluator）
  \State 按 $K_2=(E,n_{\mathrm{spill}},T)$ 或 $K_3=(T,E,n_{\mathrm{spill}})$ 更新 $A^\star$
\EndFor
\State 返回 $A^\star$；实验 harness 对最终工件运行 canonical evaluator
\end{algorithmic}
\end{algorithm}

算法~\ref{alg:supp-portfolio-zh} 最多考察
$4\times2\times|\mathcal H|$ 个组合成员；放置失败者被跳过。它对成功候选
生成地址/溢出工件并计算真实目标，不用 live-area 或 overflow-area 代理量
作最终选优。生产 \texttt{solve} 的候选循环本身不逐个调用 canonical
evaluator；实验 harness 对最终选中工件作标准验证，本文只报告零违规结果。表
~\ref{tab:supp-implementation-constants-zh} 给出正文伪代码省略但会影响
复现的确定性规则。

\begin{table}[H]
\centering
\caption{Canonical 求解器的固定常数与边界规则。}
\label{tab:supp-implementation-constants-zh}
\begin{tabular}{@{}p{.22\linewidth}p{.15\linewidth}p{.55\linewidth}@{}}
\toprule
机制 & 数值 & 精确作用\\
\midrule
前沿平局裁决 & FREE 最大 ID；操作最小 ID & ALLOC 键为（最小可解锁消费者 ID，ALLOC ID）；无可解锁消费者时第一项取无穷。\\
下次使用截断 & $2^{20}$ & 形成任一 victim 评分前，截断按基础序位置计算的距离。\\
尺寸异质性 & 最大/最小 $\ge8$ & 若该存储类型完全没有 backed 缓冲，自适应策略优先释放大连续段；否则保持距离优先。\\
backed 占比 & $\le25\%$ & 有可驱逐 backed victim 且其字节占可驱逐驻留字节不超过阈值时，优先使用这一低价储备。\\
提前重载 & 距离 $\le H$ & 仅使用已经存在的合适空洞；预取不主动制造 victim。\\
pinned 搬移 & 每请求至多 100 次 & 只有碎片化且无非 pinned victim 时，溢出并立即换入最小 pinned 缓冲；超过上限则该候选失败。\\
\bottomrule
\end{tabular}
\end{table}

\subsection{一个可逐步复算的小图}

表~\ref{tab:supp-toy-trace-zh} 给出一个按当前 best-fit/spill 实现与
canonical evaluator 规则手工展开的教学例子，用来明确溢出节点、地址与
计费如何共同形成 P2 工件；它不是运行日志、公开基准或 oracle 数据，也不
用于论文数值主张。把七个原始节点按下列顺序编号，并以相邻节点之间的六条
边 $0\to1\to2\to3\to4\to5\to6$ 构成一条 DAG 链；因此原始拓扑序可直接
复核且唯一。设某
存储容量为 4，缓冲区 $a$ 为 2 字节 backed 输入，$b$ 为 3 字节 generated
中间量。原始合法序为
\[
 \operatorname{ALLOC}(a),\ \operatorname{COPYIN}(a),\
 \operatorname{ALLOC}(b),\ \operatorname{PRODUCE}(b),\
 \operatorname{FREE}(b),\ \operatorname{USE}(a),\ \operatorname{FREE}(a).
\]
$a$ 在 COPY-IN 后到 USE 前没有强制访问。分配 $b$ 时总逻辑活跃量为 5，
best-fit 无法找到 3 字节空洞；选择换出 $a$ 后，$b$ 可占 $[0,3)$。由于
FREE 是强制事件，$b$ 必须保持驻留到其 FREE；随后再把 $a$ 换入并执行
USE。这里选择 backed 的 $a$，一次中断只收费 2 字节。

\begin{table}[H]
\centering
\caption{容量为 4 的教学小图：从原始序到增广物理工件。$o_a,i_a$ 为合成换出/换入节点。}
\label{tab:supp-toy-trace-zh}
\begin{tabular}{@{}clllr@{}}
\toprule
步 & 执行事件 & 事件后驻留段 & 空闲区间 & 累计 $E$\\
\midrule
1 & $\operatorname{ALLOC}(a)$ & $a:[0,2)$ & $[2,4)$ & 0\\
2 & $\operatorname{COPYIN}(a)$ & $a:[0,2)$ & $[2,4)$ & 0\\
3 & $o_a$ & 无 & $[0,4)$ & 2\\
4 & $\operatorname{ALLOC}(b)$ & $b:[0,3)$ & $[3,4)$ & 2\\
5 & $\operatorname{PRODUCE}(b)$ & $b:[0,3)$ & $[3,4)$ & 2\\
6 & $\operatorname{FREE}(b)$ & 无 & $[0,4)$ & 2\\
7 & $i_a$ & $a:[0,2)$ & $[2,4)$ & 2\\
8 & $\operatorname{USE}(a)$ & $a:[0,2)$ & $[2,4)$ & 2\\
9 & $\operatorname{FREE}(a)$ & 无 & $[0,4)$ & 2\\
\bottomrule
\end{tabular}
\end{table}

该工件的溢出表只有 $(a,0)$ 一项，所以
$n_{\mathrm{spill}}=1$、$C=2,D=0,V=2,E=2$。验证器还会检查原始依赖、
$\operatorname{ALLOC}(a)\prec o_a\prec i_a\prec\operatorname{FREE}(a)$、
$o_a$ 与 $i_a$ 之间无 $a$ 的使用，以及所有半开地址区间均不越界或时空
重叠。若换出的是同尺寸 generated 缓冲，流量会是 4；这正是命题 1 的
单位代价差异，而非“排序自动使用类别”的证据。

\subsection{真实实例 Conv\_Case0 的端到端路径}

真实实例不能逐节点印出 2,580 行，表~\ref{tab:supp-conv0-path-zh} 因而按
求解角色展示从候选序到验证工件的完整方法路径。可扩展 v2 最多考察 24 个
P2 组合成员，逐一执行“排序--best-fit--victim--溢出插入--内部计价/计时”，
再由实验 harness 对最终选中工件作 canonical 验证。该工件零违规，
$E=66{,}828$、140 次溢出；其流量可直接分解为
$C=12{,}672,D=27{,}078,V=39{,}750$，故
$E=V+D=39{,}750+27{,}078=66{,}828$。

\begin{table}[H]
\centering
\caption{Conv\_Case0 的方法路径。exact 行只证固定 \texttt{unlock-frontier} 序的流量最优。}
\label{tab:supp-conv0-path-zh}
\begin{tabular}{@{}llrrp{.34\linewidth}@{}}
\toprule
角色 & 工件/搜索范围 & $E$ & 溢出数 & 可作出的结论\\
\midrule
历史前身 & iter038 & 88,044 & 339 & 可扩展 P2 前身\\
生产 & scalable v2，24 候选 & 66,828 & 140 & canonical-valid 生产结果\\
离线修复 & 10,000 次随机单节点 proposal & 65,532 & 137 & Conv0 个案，非默认阶段\\
固定序 exact & frontier gap+packing & 57,408 & 112 & 下界相等、打包成功、固定序流量证书\\
官方 & 官方 P2 工件 & 73,500 & 141 & 同口径比较对象\\
\bottomrule
\end{tabular}
\end{table}

固定序 exact 行进一步展示证书链：间隙 CP-SAT 在 0.88 秒的记录运行中给出
下界与目标均为 57,408；连续地址打包状态为 OPTIMAL（这里只表示可满足性），
标准 schedule/memory/spill 文件生成成功，canonical evaluator 返回零违规
且同样计算出 57,408。由定理 3，它证明该 \texttt{unlock-frontier} 序的
最小\emph{流量}；它没有证明 112 是同流量下最少溢出数，也没有搜索所有
拓扑序。作为排序敏感性的受控证据，同一 planner 固定为 P1 序时证明的流量
最优值为 81,504；这一对比只对 Conv0 成立。

\subsection{离线代价修复与固定序 oracle}

离线修复脚本观察高代价溢出前沿，提出保持依赖合法的节点或依赖链移动，
重新运行物理规划，并只接受 P2 键严格改善的 proposal。实验预算并不统一：
Conv\_Case0 使用 10,000 次随机单节点 hill search；Conv\_Case1 使用 beam
宽度 10、两轮 frontier search；两个 FA 实例各尝试 2,000 个随机 proposal
但没有接受移动；两个 Matmul 实例没有运行修复。因而它只能说明某些个案中
代价语义可能影响排序，不能称为六例统一算法或普遍改进。

固定序 oracle 先解式~\eqref{eq:supp-gap-obj-zh} 的加权间隙选择，再进行
若干确定性贪心打包；必要时调用独立 \texttt{NoOverlap2D} 可满足性模型。
最后它输出与生产求解器相同格式的 schedule/memory/spill 工件，并调用同一
evaluator。8 个内部 17 节点 oracle 实例上，生产求解器在其\emph{自身所选
固定序}的流量依次为 352、1,408、704、704、352、1,408、704 和 704；
每一项都与该固定序的 CP-SAT 下界相等且通过验证。这是 8 个离散小图的
固定序吻合，不是大规模扩展性或跨序全局最优性证据。

\subsection{失败、未完成与复现边界}

完整求解过程也包括失败状态，而不是只保留成功证书：
\begin{itemize}
  \item FA\_Case1 的固定序工件为 canonical-valid FEASIBLE：$E=32{,}512$，
        下界 31,936，gap 为 576；不能称为 certified optimal；
  \item Matmul\_Case0 的可复现 120.09 秒运行给出 FEASIBLE 34,816、下界
        29,952；同一序另有合法启发式工件 34,688，故审计区间为
        $[29{,}952,34{,}688]$，旧的“34,688 OPTIMAL”说法已撤销；
  \item Conv\_Case1 未在研究预算内完成 exact+packing，Matmul\_Case1 未运行；
  \item P3 在六个公开实例中对官方时间是 5 胜 1 负，Conv\_Case1 为
        1,118,687 对 1,073,322，不能写成普遍占优；
  \item 汇总 JSON 可核对指标和验证状态，但若其旁边没有配套文本工件，独立
        replay 需要重跑确定性求解器。固定序证书则以机器 JSON 和匹配的
        schedule/memory/spill 文件为准，旧的手写状态不得覆盖机器记录。
\end{itemize}

最后，生产路径的地址放置和驻留扫描在固定 portfolio 下是可扩展实现；
CP-SAT 的间隙选择与后备 packing 最坏情况下为指数复杂度。当前
canonical \texttt{solve} 没有 CP-SAT 规模阈值、超时切换或回退集成，故
不能把 oracle 的成功结果描述为默认求解器能力。
